% GENERATED FILE. Edit canonical lessons, then run npm run generate:books.
% canonical-source-hash: fnv1a64:f78d29c0b4318e29
% canonical-lessons: TE-C417-nibandhana, TE-C417-arhata, TE-C417-anumatincu, TE-C417-minaha, TE-C417-vartincu

\chapter{Rules and Requirements}
\label{ch:te-rules-and-requirements}
\hlchaptermodality{417}

\begin{chapteropening}
\textbf{By the end of this chapter:} \emph{I can talk about clauses, qualifications and what a rule allows.}

Complete the last lesson of chapter 417: I can talk about clauses, qualifications and what a rule allows.
\end{chapteropening}

\section[\texorpdfstring{\te{నిబంధన} (nibandhana)}{నిబంధన (nibandhana)}]{\te{నిబంధన} (nibandhana) — a stipulation; a clause}

\label{lesson:TE-C417-nibandhana}



\begin{quote}
Before the new one: say the Telugu for a syllabus topic, then the Telugu for to review.
\end{quote}

\subsection*{You'll want to know: \te{నిబంధన}}
\textbf{\te{నిబంధన}} — \emph{nibandhana} — \textquotedblleft{}a stipulation; a clause\textquotedblright{}.

\begin{quote}
\te{ఒప్పందంలో} \te{ఒక} \te{కొత్త} \te{నిబంధన} \te{చేర్చారు}.
\end{quote}

\begin{quote}
\emph{oppandaṁlō oka kotta nibandhana cērcāru.}
\end{quote}

\begin{quote}
They added a new clause to the agreement.
\end{quote}

Plural: \te{నిబంధనలు}, \emph{nibandhanalu}.

Plural \te{నిబంధనలు}, terms and conditions.

\begin{cousinweb}
\textbf{\te{నిబంధన}} (\emph{nibandhana}) is borrowed from Sanskrit \emph{nibandhana}, \textquotedblleft{}binding, fastening; a band, a support\textquotedblright{} (Monier-Williams).
\end{cousinweb}

\begin{grammarlens}[title={What you've built}]
One more word for saying what happens if.
\end{grammarlens}

\subsection*{Your turn}
\begin{itemize}
\raggedright
  \item \emph{Say it:} \emph{nibandhana}
  \item \emph{Say it:} \emph{nibandhana}, once more
  \item \emph{Say it:} \emph{oppandaṁlō oka kotta nibandhana cērcāru.} — They added a new clause to the agreement.
  \item \emph{Recall:} say the Telugu for a syllabus topic, then the Telugu for to review, then say \emph{nibandhana} again
\end{itemize}

\begin{tcolorbox}[breakable,colback=teal!4,colframe=teal!35!black,title={Before you move on}]
What does \te{నిబంధన} mean? (\textquotedblleft{}A stipulation; a clause\textquotedblright{}.) Say it once more.
\end{tcolorbox}
\section[\texorpdfstring{\te{అర్హత} (arhata)}{అర్హత (arhata)}]{\te{అర్హత} (arhata) — eligibility; a qualification}

\label{lesson:TE-C417-arhata}



\begin{quote}
Before the new one: say the Telugu for to review, then the Telugu for a stipulation.
\end{quote}

\subsection*{You'll want to know: \te{అర్హత}}
\textbf{\te{అర్హత}} — \emph{arhata} — \textquotedblleft{}eligibility; a qualification\textquotedblright{}.

\begin{quote}
\te{ఆ} \te{ఉద్యోగానికి} \te{నాకు} \te{అర్హత} \te{ఉంది}.
\end{quote}

\begin{quote}
\emph{ā udyōgāniki nāku arhata undi.}
\end{quote}

\begin{quote}
I am qualified for that job.
\end{quote}

With the -\te{తే} of a condition: \te{అర్హత} \te{ఉంటే}, if you qualify.

\begin{cousinweb}
\textbf{\te{అర్హత}} (\emph{arhata}) is built on Sanskrit \emph{arha}, \textquotedblleft{}meriting, deserving, worthy of\textquotedblright{} (Monier-Williams), with the abstract ending -\te{త}.
\end{cousinweb}

\begin{grammarlens}[title={What you've built}]
One more word for saying what happens if.
\end{grammarlens}

\subsection*{Your turn}
\begin{itemize}
\raggedright
  \item \emph{Say it:} \emph{arhata}
  \item \emph{Say it:} \emph{arhata}, once more
  \item \emph{Say it:} \emph{ā udyōgāniki nāku arhata undi.} — I am qualified for that job.
  \item \emph{Recall:} say the Telugu for to review, then the Telugu for a stipulation, then say \emph{arhata} again
\end{itemize}

\begin{tcolorbox}[breakable,colback=teal!4,colframe=teal!35!black,title={Before you move on}]
What does \te{అర్హత} mean? (\textquotedblleft{}Eligibility; a qualification\textquotedblright{}.) Say it once more.
\end{tcolorbox}
\section[\texorpdfstring{\te{అనుమతించు} (anumatiñcu)}{అనుమతించు (anumatiñcu)}]{\te{అనుమతించు} (anumatiñcu) — to permit; to allow}

\label{lesson:TE-C417-anumatincu}



\begin{quote}
Before the new one: say the Telugu for a stipulation, then the Telugu for eligibility.
\end{quote}

\subsection*{You'll want to know: \te{అనుమతించు}}
\textbf{\te{అనుమతించు}} — \emph{anumatiñcu} — \textquotedblleft{}to permit; to allow\textquotedblright{}.

\begin{quote}
\te{మా} \te{నాన్న} \te{నన్ను} \te{రాత్రి} \te{బయటకు} \te{వెళ్ళడానికి} \te{అనుమతించాడు}.
\end{quote}

\begin{quote}
\emph{mā nānna nannu rātri bayaṭaku veḷḷaḍāniki anumatiñcāḍu.}
\end{quote}

\begin{quote}
My father allowed me to go out at night.
\end{quote}

Past (I …): \te{అనుమతించాను}, \emph{anumatiñcānu}.

The person takes -\te{ని} and the act -\te{డానికి}. The noun \te{అనుమతి}, permission, you know.

\begin{cousinweb}
\textbf{\te{అనుమతించు}} (\emph{anumatiñcu}) is built on Sanskrit \emph{anumati}, \textquotedblleft{}assent, permission, approbation\textquotedblright{} (Monier-Williams), with the Telugu verb ending -\te{ఇంచు}; \te{అనుమతి} is the same word.
\end{cousinweb}

\begin{grammarlens}[title={What you've built}]
One more word for saying what happens if.
\end{grammarlens}

\subsection*{Your turn}
\begin{itemize}
\raggedright
  \item \emph{Say it:} \emph{anumatiñcu}
  \item \emph{Say it:} \emph{anumatiñcu}, once more
  \item \emph{Say it:} \emph{mā nānna nannu rātri bayaṭaku veḷḷaḍāniki anumatiñcāḍu.} — My father allowed me to go out at night.
  \item \emph{Recall:} say the Telugu for a stipulation, then the Telugu for eligibility, then say \emph{anumatiñcu} again
\end{itemize}

\begin{tcolorbox}[breakable,colback=teal!4,colframe=teal!35!black,title={Before you move on}]
What does \te{అనుమతించు} mean? (\textquotedblleft{}To permit; to allow\textquotedblright{}.) Say it once more.
\end{tcolorbox}
\section[\texorpdfstring{\te{మినహా} (minahā)}{మినహా (minahā)}]{\te{మినహా} (minahā) — except for; apart from}

\label{lesson:TE-C417-minaha}



\begin{quote}
Before the new one: say the Telugu for eligibility, then the Telugu for to permit.
\end{quote}

\subsection*{You'll want to know: \te{మినహా}}
\textbf{\te{మినహా}} — \emph{minahā} — \textquotedblleft{}except for; apart from\textquotedblright{}.

\begin{quote}
\te{ఆదివారం} \te{మినహా} \te{ప్రతి} \te{రోజు} \te{దుకాణం} \te{తెరిచి} \te{ఉంటుంది}.
\end{quote}

\begin{quote}
\emph{ādivāraṁ minahā prati rōju dukāṇaṁ terici uṇṭundi.}
\end{quote}

\begin{quote}
The shop is open every day except Sunday.
\end{quote}

It follows the noun it leaves out: \te{ఆదివారం} \te{మినహా}, except Sunday.

\begin{cousinweb}
\textbf{\te{మినహా}} (\emph{minahā}) is usually traced to Hindustani \emph{minhā}, \textquotedblleft{}deducted, left out\textquotedblright{}, a word of Arabic origin.
\end{cousinweb}

\begin{grammarlens}[title={What you've built}]
One more word for saying what happens if.
\end{grammarlens}

\subsection*{Your turn}
\begin{itemize}
\raggedright
  \item \emph{Say it:} \emph{minahā}
  \item \emph{Say it:} \emph{minahā}, once more
  \item \emph{Say it:} \emph{ādivāraṁ minahā prati rōju dukāṇaṁ terici uṇṭundi.} — The shop is open every day except Sunday.
  \item \emph{Recall:} say the Telugu for eligibility, then the Telugu for to permit, then say \emph{minahā} again
\end{itemize}

\begin{tcolorbox}[breakable,colback=teal!4,colframe=teal!35!black,title={Before you move on}]
What does \te{మినహా} mean? (\textquotedblleft{}Except for; apart from\textquotedblright{}.) Say it once more.
\end{tcolorbox}
\section[\texorpdfstring{\te{వర్తించు} (vartiñcu)}{వర్తించు (vartiñcu)}]{\te{వర్తించు} (vartiñcu) — to apply (to); to be relevant}

\label{lesson:TE-C417-vartincu}



\begin{quote}
Before the new one: say the Telugu for to permit, then the Telugu for except for.
\end{quote}

\subsection*{You'll want to know: \te{వర్తించు}}
\textbf{\te{వర్తించు}} — \emph{vartiñcu} — \textquotedblleft{}to apply (to); to be relevant\textquotedblright{}.

\begin{quote}
\te{ఆ} \te{నియమం} \te{పిల్లలకు} \te{వర్తించదు}.
\end{quote}

\begin{quote}
\emph{ā niyamaṁ pillalaku vartiñcadu.}
\end{quote}

\begin{quote}
That rule does not apply to children.
\end{quote}

Past (it …): \te{వర్తించింది}, \emph{vartiñcindi}.

What it applies to takes -\te{కి}. Rules apply, not people, so its past is shown as \te{వర్తించింది}, it applied.

\begin{cousinweb}
\textbf{\te{వర్తించు}} (\emph{vartiñcu}) is built on Sanskrit \emph{vartana}, \textquotedblleft{}turning, rolling\textquotedblright{} (CDIAL 11354), and also \textquotedblleft{}abiding, staying\textquotedblright{} (Monier-Williams), with the Telugu verb ending -\te{ఇంచు}.
\end{cousinweb}

\begin{grammarlens}[title={What you've built}]
One more word for saying what happens if.
\end{grammarlens}

\subsection*{Your turn}
\begin{itemize}
\raggedright
  \item \emph{Say it:} \emph{vartiñcu}
  \item \emph{Say it:} \emph{vartiñcu}, once more
  \item \emph{Say it:} \emph{ā niyamaṁ pillalaku vartiñcadu.} — That rule does not apply to children.
  \item \emph{Recall:} say the Telugu for to permit, then the Telugu for except for, then say \emph{vartiñcu} again
\end{itemize}

\begin{tcolorbox}[breakable,colback=teal!4,colframe=teal!35!black,title={Before you move on}]
What does \te{వర్తించు} mean? (\textquotedblleft{}To apply (to); to be relevant\textquotedblright{}.) Say it once more.
\end{tcolorbox}
